\chapter*{Ажлын төлөвлөгөө}
\addcontentsline{toc}{chapter}{Ажлын төлөвлөгөө}
\label{ch:plan}

\noindent\textbf{Оюутны нэр:} \authorShort \hfill \textbf{Удирдагч багшийн нэр:} \supervisorname

\noindent\textbf{Судалгааны ажлын сэдэв:} \titlename % сэдвийн нэрийг удирдагчтай баталгаажуулна

\vspace{0.2cm}

Ажлыг долоо хоног бүрээр төлөвлөсөн. Долоо хоног бүр нэг гол ажил хийж, тухайн ажилтай холбоотой нэг судалгааны өгүүлэл уншина. Улаанаар тэмдэглэсэн нь сургуулийн гарын авлагад заасан үзлэг, хамгаалалт юм.

{\small
\setlength{\tabcolsep}{4pt}
\renewcommand{\arraystretch}{1.12}
\begin{longtable}{|p{1.0cm}|p{0.9cm}|p{1.85cm}|p{6.0cm}|p{3.0cm}|p{1.3cm}|}
\caption{Ажлын долоо хоногийн төлөвлөгөө (2026.10--2027.05)}\label{tab:plan}\\
\hline
\textbf{Сар} & \textbf{7 хоног} & \textbf{Огноо} & \textbf{Гол ажил} & \textbf{Судалгаа (унших)} & \textbf{Гүйцэт\-гэл} \\
\hline
\endfirsthead
\hline
\textbf{Сар} & \textbf{7 хоног} & \textbf{Огноо} & \textbf{Гол ажил} & \textbf{Судалгаа (унших)} & \textbf{Гүйцэт\-гэл} \\
\hline
\endhead
10 сар & W01 & 10.26--11.01 & Ажлаа цэгцэлж, системийн гол сонголтуудыг хийх & -- & $\checkmark$ \\
\hline
11 сар & W02 & 11.02--11.08 & Семинар хийж, багштай тодруулах; Whisper--Gemini харьцуулалтыг дипломд бичих & RAG & Явж байна \\
 & W03 & 11.09--11.15 & Дуу таних аргаа шинэ бичлэгээр харьцуулж сонгох & Хариугүй асуулт (SQuAD 2.0) & \\
 & W04 & 11.16--11.22 & Шалгах асуултын багц бэлдэж түгжих & BM25 & \\
 & W05 & 11.23--11.29 & Үзлэг №1-ийн материал бэлдэх & RRF, тойм бичих & \\
\hline
12 сар & W06 & 11.30--12.06 & \textcolor{red}{Үзлэг №1 (12/01)} & Давтлага & \\
 & W07 & 12.07--12.13 & Хайлтын аргуудыг харьцуулах (мэдлэгийн модтой, модгүй) & Multilingual E5 & \\
 & W08 & 12.14--12.20 & Хариулт сонгох хоёр аргыг (дүрэм, LLM) харьцуулах & Эшлэлтэй хариулт & \\
 & W09 & 12.21--12.27 & Мэдлэгийн модыг баяжуулах: багш, өрөө, анги, хуваарь & RAPTOR, GraphRAG & \\
 & W10 & 12.28--01.03 & Завсрын дүгнэлт хийх & Давтлага & \\
\hline
1 сар & W11 & 01.04--01.10 & Нэр, тоо, оныг дуугаар зөв уншуулах & Тоо хэвийн болгох & \\
 & W12 & 01.11--01.17 & Хариултыг сонсож ойлгосныг шалгах & TTS-ийн ойлгомж & \\
 & W13 & 01.18--01.24 & Хариулт ба дуу танилтын гол алдааг засах & Дуу танилтын алдаа ба нэр & \\
 & W14 & 01.25--01.31 & Мэдээлэл шинэчлэх, тасалдлыг шалгах & DPR & \\
\hline
2 сар & W15 & 02.01--02.07 & Бичсэн ба ярьсан асуултыг харьцуулах; \textcolor{red}{Үзлэг №2} & Утгын хэмжүүр & \\
 & W16 & 02.08--02.14 & Өөр компьютер дээр ажиллуулах & Хэрэглэгчийн хүлээлт & \\
 & W17 & 02.15--02.21 & Эцсийн хувилбар, LLM-ийн эцсийн шийдвэр & LLM ба RAG & \\
 & W18 & 02.22--02.28 & Гол хэмжилтүүдийг хийх & RAGAS & \\
\hline
3 сар & W19 & 03.01--03.07 & Хүний үнэлгээ хийх & TTS-ийн хүний үнэлгээ & \\
 & W20 & 03.08--03.14 & Үр дүн, алдааны шинжилгээг бичих & Bootstrap статистик & \\
 & W21 & 03.15--03.21 & \textcolor{red}{Бүтэн ноорог (40 хүртэл хуудас)} & Давтлага & \\
 & W22 & 03.22--03.28 & Удирдагчийн саналаар засах & -- & \\
 & W23 & 03.29--04.04 & Үзүүлэн, хамгаалалтын яриа бэлдэх & -- & \\
\hline
4 сар & W24 & 04.05--04.11 & Урьдчилсан хамгаалалтын материал & -- & \\
 & W25 & 04.12--04.18 & \textcolor{red}{Урьдчилсан хамгаалалт} & -- & \\
 & W26 & 04.19--04.25 & Комиссын саналаар засах & -- & \\
 & W27 & 04.26--05.02 & Засварыг дуусгах, давтан ажиллуулах заавар & -- & \\
\hline
5 сар & W28 & 05.03--05.09 & Эцсийн бичвэр, Монгол ба англи хураангуй & -- & \\
 & W29 & 05.10--05.16 & Товчлол (A4, 2--4 хуудас), хүлээлгэн өгөх материал & -- & \\
 & W30 & 05.17--05.23 & Хамгаалалтын давтлага & -- & \\
 & W31 & 05.24--05.30 & \textcolor{red}{Эцсийн хяналт, хүлээлгэн өгөх} & -- & \\
 & W32 & 05.31 & Нөөц өдөр & -- & \\
\hline
\end{longtable}
}

% Үзлэг №2, урьдчилсан ба эцсийн хамгаалалтын яг огноог сургууль зарлахаар шинэчилнэ.
